% GENERATED by scripts/generate-mindmaps.mjs — do not hand-edit.
\begin{tikzpicture}[
  mindmap,
  every node/.style={
    concept,
    execute at begin node=\hskip0pt,
    align=center,
    font=\sffamily
  },
  mm root/.style={
    concept color=graphite,
    text=ghostwhite,
    font=\sffamily\bfseries,
    minimum size=3.35cm,
    text width=2.75cm
  },
  mm section/.style={
    font=\sffamily\bfseries\scriptsize,
    minimum size=2.42cm,
    text width=1.92cm
  },
  mm topic/.style={
    font=\sffamily\tiny,
    minimum size=1.92cm,
    text width=1.46cm
  }
]
  \node[mm root] (root) at (0,0) {Chapter 3\\Automata over Infinite Words};
  \node[mm section, concept color=brightsky] (s1) at (82.17:5.1cm) {Infinite Words};
  \node[mm topic, concept color=brightsky!44!white] (t1_1) at (82.17:9.8cm) {\ensuremath{\omega}-word};
  \node[mm section, concept color=vvteal] (s2) at (35.22:5.1cm) {From Finite Languages to Infinite Languages};
  \node[mm topic, concept color=vvteal!44!white] (t2_1) at (66.52:9.8cm) {\ensuremath{\omega}-power};
  \node[mm topic, concept color=vvteal!44!white] (t2_2) at (50.87:9.8cm) {Segments and suffixes};
  \node[mm topic, concept color=vvteal!44!white] (t2_3) at (35.22:9.8cm) {Infinitely/finitely many positions};
  \node[mm topic, concept color=vvteal!44!white] (t2_4) at (19.57:9.8cm) {Vectorial closure};
  \node[mm topic, concept color=vvteal!44!white] (t2_5) at (3.91:9.8cm) {Vectorial closure on the running example};
  \node[mm section, concept color=vvviolet] (s3) at (-66.52:5.1cm) {Büchi Automata};
  \node[mm topic, concept color=vvviolet!44!white] (t3_1) at (-11.74:9.8cm) {Büchi automaton};
  \node[mm topic, concept color=vvviolet!44!white] (t3_2) at (-27.39:9.8cm) {Infinitely often};
  \node[mm topic, concept color=vvviolet!44!white] (t3_3) at (-43.04:9.8cm) {Infinity sets in the running example};
  \node[mm topic, concept color=vvviolet!44!white] (t3_4) at (-58.7:9.8cm) {Run and Büchi acceptance};
  \node[mm topic, concept color=vvviolet!44!white] (t3_5) at (-74.35:9.8cm) {\ensuremath{\omega}-regular language};
  \node[mm topic, concept color=vvviolet!44!white] (t3_6) at (-90:9.8cm) {Infinitely many saves};
  \node[mm topic, concept color=vvviolet!44!white] (t3_7) at (-105.65:9.8cm) {Parity between occurrences};
  \node[mm topic, concept color=vvviolet!44!white] (t3_8) at (-121.3:9.8cm) {Complement of a Büchi language};
  \node[mm section, concept color=vvorange] (s4) at (-168.26:5.1cm) {Closure and Expressions};
  \node[mm topic, concept color=vvorange!44!white] (t4_1) at (-136.96:9.8cm) {Basic closure properties};
  \node[mm topic, concept color=vvorange!44!white] (t4_2) at (-152.61:9.8cm) {Constructions};
  \node[mm topic, concept color=vvorange!44!white] (t4_3) at (-168.26:9.8cm) {Vectorial closure of a regular language};
  \node[mm topic, concept color=vvorange!44!white] (t4_4) at (-183.91:9.8cm) {\ensuremath{\omega}-regular expression};
  \node[mm topic, concept color=vvorange!44!white] (t4_5) at (-199.57:9.8cm) {Büchi-expression equivalence};
  \node[mm section, concept color=tomato] (s5) at (-238.7:5.1cm) {Emptiness and Lassos};
  \node[mm topic, concept color=tomato!44!white] (t5_1) at (-215.22:9.8cm) {Decidability of Büchi emptiness};
  \node[mm topic, concept color=tomato!44!white] (t5_2) at (-230.87:9.8cm) {Ultimately periodic word};
  \node[mm topic, concept color=tomato!44!white] (t5_3) at (-246.52:9.8cm) {Ultimately periodic witnesses};
  \node[mm topic, concept color=tomato!44!white] (t5_4) at (-262.17:9.8cm) {Languages without periodic words};

  \begin{scope}[on background layer, line cap=round]
    \draw[brightsky, line width=11pt] (root) -- (s1);
    \draw[brightsky!60!white, line width=5.5pt] (s1) -- (t1_1);
    \draw[vvteal, line width=11pt] (root) -- (s2);
    \draw[vvteal!60!white, line width=5.5pt] (s2) -- (t2_1);
    \draw[vvteal!60!white, line width=5.5pt] (s2) -- (t2_2);
    \draw[vvteal!60!white, line width=5.5pt] (s2) -- (t2_3);
    \draw[vvteal!60!white, line width=5.5pt] (s2) -- (t2_4);
    \draw[vvteal!60!white, line width=5.5pt] (s2) -- (t2_5);
    \draw[vvviolet, line width=11pt] (root) -- (s3);
    \draw[vvviolet!60!white, line width=5.5pt] (s3) -- (t3_1);
    \draw[vvviolet!60!white, line width=5.5pt] (s3) -- (t3_2);
    \draw[vvviolet!60!white, line width=5.5pt] (s3) -- (t3_3);
    \draw[vvviolet!60!white, line width=5.5pt] (s3) -- (t3_4);
    \draw[vvviolet!60!white, line width=5.5pt] (s3) -- (t3_5);
    \draw[vvviolet!60!white, line width=5.5pt] (s3) -- (t3_6);
    \draw[vvviolet!60!white, line width=5.5pt] (s3) -- (t3_7);
    \draw[vvviolet!60!white, line width=5.5pt] (s3) -- (t3_8);
    \draw[vvorange, line width=11pt] (root) -- (s4);
    \draw[vvorange!60!white, line width=5.5pt] (s4) -- (t4_1);
    \draw[vvorange!60!white, line width=5.5pt] (s4) -- (t4_2);
    \draw[vvorange!60!white, line width=5.5pt] (s4) -- (t4_3);
    \draw[vvorange!60!white, line width=5.5pt] (s4) -- (t4_4);
    \draw[vvorange!60!white, line width=5.5pt] (s4) -- (t4_5);
    \draw[tomato, line width=11pt] (root) -- (s5);
    \draw[tomato!60!white, line width=5.5pt] (s5) -- (t5_1);
    \draw[tomato!60!white, line width=5.5pt] (s5) -- (t5_2);
    \draw[tomato!60!white, line width=5.5pt] (s5) -- (t5_3);
    \draw[tomato!60!white, line width=5.5pt] (s5) -- (t5_4);
  \end{scope}
\end{tikzpicture}
